\centerline{\bf \Huge Солянка I}
\centerline{\bf \Huge }

\begin{flushright}
        \scriptsize{}
\end{flushright}

\problem{1}
В группе ЭлМШ 20 человек. Из них все, кроме Арслана дружат ровно с 5 одноклассниками. Может ли Арслан ни с кем не дружить?


\problem{2}
10 учеников решали дз Расула Владимировича из 10 задач. Выяснилось, что каждый ученик решил ровно две задачи, и каждую задачу решило два ученика. 
Докажите, что Расул Владимирович может устроить разбор дз таким образом, чтобы каждый ученик рассказал решение ровно одной задачи и все задачи были бы разобраны. Ученик может разбирать только ту задачу, которую он решил.

\problem{3}
Сколько существует девятизначных чисел, у которых сумма 

(a) последних двух;

(b) всех; цифр четна?


\problem{4}
Сколькими способами можно поставить на шахматную доску двух королей так, чтобы они не били друг друга?  

\problem{5}
Сколькими способами можно поставить на шахматную доску двух королей так, чтобы они не били друг друга?